%% corollary3.tex — Selective Consolidation
%% Fixes Issue: #23 (definitional gap for "invariant coordinates")
\section{Corollary 3: Selective Consolidation}
\label{sec:cor3}

%% ======================================================================
\subsection{Formal Definition of Invariant Coordinates}
%% ======================================================================

\begin{definition}[Gradient-sign stability --- gauge-invariant formulation]
\label{def:gradient-sign-stability}
We work with the \emph{projection-matrix gradient}
$G^{(m)} = \nabla_P \hat{L}_{mB}(P) \in \R^{d \times d}$ (symmetric),
evaluated at $P = \hat{U}\hat{U}^\top$. This is gauge-invariant:
$G^{(m)}$ depends on the subspace $\col(\hat{U})$, not on the
choice of orthonormal basis.

For each entry $(i,j)$ with $i \leq j$ (upper triangle),
define the \emph{entry-wise sign stability} over $W$ consecutive
consolidation rounds:
\[
  s_{ij} \coloneqq \frac{1}{W}\left|\sum_{m=1}^W
  \sign\bigl(G^{(m)}_{ij}\bigr)\right| \in [0, 1].
\]
An entry $(i,j)$ is \emph{stable} if $s_{ij} \geq \tau_s$ for a
threshold $\tau_s \in (0, 1]$, and \emph{transient} otherwise.
The stable index set is $\mathcal{I} = \{(i,j) : s_{ij} \geq \tau_s\}$.
\end{definition}

\begin{remark}[Interpretation]
A coordinate with $s_j \approx 1$ has gradient consistently pointing
in one direction across consolidation rounds, indicating a persistent
signal. A coordinate with $s_j \approx 0$ has gradient sign flipping
randomly, indicating noise or transient effects.
\end{remark}

%% ======================================================================
\subsection{Selective Consolidation Guarantee}
%% ======================================================================

\begin{definition}[Selective consolidation]
\label{def:selective-consolidation}
Given the stable set $\mathcal{I}$, \emph{selective consolidation}
applies the projection-matrix gradient update only to stable entries:
\[
  \hat{P}^{\mathrm{new}}_{ij} =
  \begin{cases}
    \hat{P}^{\mathrm{old}}_{ij} - \eta_s G_{ij}
      & \text{if } (i,j) \in \mathcal{I}, \\
    \hat{P}^{\mathrm{old}}_{ij}
      & \text{otherwise},
  \end{cases}
\]
followed by a retraction step $\hat{P}^{\mathrm{new}} \gets
\mathrm{Proj}_{\mathcal{P}_r}(\hat{P}^{\mathrm{new}})$ (eigenvalue
thresholding to the nearest rank-$r$ projector) to restore feasibility.
\emph{Indiscriminate consolidation} applies the gradient update to all
entries before retraction.
\end{definition}

\begin{proposition}[Selective consolidation bound --- conditional]
\label{prop:selective}
For each entry $(i,j)$ of the projection-matrix gradient, decompose the
stochastic gradient as $G_{ij} = \mu_{ij} + \xi_{ij}$, where
$\mu_{ij} = \E[G_{ij}]$ is the signal and $\xi_{ij}$ is zero-mean
noise with $\E[\xi_{ij}^2] = \sigma_{ij}^2$. Assume the gradient
entries $\{G_{ij}^{(m)}\}_{m=1}^W$ are independent across
consolidation rounds (valid when different data batches are used).

\textbf{Sign-stability characterization:}
\begin{enumerate}[nosep]
  \item The expected sign-stability satisfies
        $\E[s_{ij}] = |2p_{ij} - 1| + O(W^{-1/2})$,
        where $p_{ij} = \Pr(G_{ij} > 0)$. Under Gaussian noise:
        $p_{ij} = \Phi(\mu_{ij}/\sigma_{ij})$ and
        $|2\Phi(x)-1| \leq \sqrt{2/\pi}\,|x|$ for $|x| \leq 1$.
  \item By Hoeffding's inequality applied to the i.i.d.\ $\pm 1$ signs
        $X_m = \sign(G_{ij}^{(m)})$:
        $\Pr(s_{ij} \geq \tau_s) \leq 2\exp(-W(\tau_s - |2p_{ij}-1|)^2/2)$
        when $\tau_s > |2p_{ij}-1|$.
\end{enumerate}

\textbf{Selective vs.\ indiscriminate regret (state-dependent):}
The per-round improvement from updating entry $(i,j)$ is
\[
  \Delta R_{ij} = 2\eta_s \mu_{ij}\, e_{ij}
  - \eta_s^2(\mu_{ij}^2 + \sigma_{ij}^2),
\]
where $e_{ij} = \hat{P}_{ij} - P_{ij}$ is the current tracking error
(with sign chosen so $\mu_{ij}\, e_{ij} \geq 0$ when the gradient
points toward the truth). Selective consolidation skips
all $(i,j) \in \mathcal{T}$. The net effect is:
\[
  \E[\reg_T^{\mathrm{selective}}] -
  \E[\reg_T^{\mathrm{indiscriminate}}]
  = \frac{T}{B}\sum_{(i,j) \in \mathcal{T}}
    \bigl(-2\eta_s |\mu_{ij}| |e_{ij}|
    + \eta_s^2(\mu_{ij}^2 + \sigma_{ij}^2)\bigr).
\]
Selective consolidation is \emph{weakly better} when, for all
$(i,j) \in \mathcal{T}$:
\[
  |\mu_{ij}|\, |e_{ij}| \leq \frac{\eta_s}{2}(\mu_{ij}^2 + \sigma_{ij}^2).
\]
This holds when either (a) $\mathrm{SNR}_{ij} \ll 1$ (noise-dominated,
regardless of $e_{ij}$), or (b) $|e_{ij}|$ is small
(near-converged subspace).
\end{proposition}

\begin{proof}
\textbf{Step 1: Sign stability.}
For each entry, $X_m = \sign(G_{ij}^{(m)})$ are i.i.d.\ (by the
independence assumption) with $\E[X_m] = 2p_{ij} - 1$ where
$p_{ij} = \Pr(G_{ij} > 0)$.
The sample mean $\bar{X} = \frac{1}{W}\sum X_m$ satisfies
$\E[\bar{X}] = 2p_{ij}-1$ and
$s_{ij} = |\bar{X}|$. By Jensen: $\E[s_{ij}] \geq |\E[\bar{X}]|
= |2p_{ij}-1|$. For the reverse direction, $\E[|\bar{X}|]
\leq |2p_{ij}-1| + \E[|\bar{X} - \E[\bar{X}]|]
\leq |2p_{ij}-1| + \sqrt{\mathrm{Var}(\bar{X})}
= |2p_{ij}-1| + O(W^{-1/2})$.

Under Gaussian noise: $p_{ij} = \Phi(\mu_{ij}/\sigma_{ij})$.
For small SNR: $|2\Phi(x)-1| \leq \sqrt{2/\pi}\,|x| < |x|$ for $|x| \leq 1$.

Hoeffding on the bounded $X_m \in \{-1,+1\}$:
$\Pr(|\bar{X}| \geq \tau_s)
\leq 2\exp(-W(\tau_s-|2p_{ij}-1|)^2/2)$ when $\tau_s > |2p_{ij}-1|$.

\textbf{Step 2: Per-entry regret from updating.}
Updating entry $(i,j)$ of the projection matrix changes the
tracking error as:
$\hat{P}_{ij}^{\mathrm{new}} = \hat{P}_{ij}^{\mathrm{old}} - \eta_s G_{ij}$.
The expected squared error change (before retraction):
\begin{align*}
  \E[(\hat{P}_{ij}^{\mathrm{new}} - P_{ij})^2]
  &= e_{ij}^2 - 2\eta_s\mu_{ij}e_{ij} + \eta_s^2(\mu_{ij}^2+\sigma_{ij}^2).
\end{align*}
The improvement $\Delta R_{ij} = 2\eta_s\mu_{ij}e_{ij}
- \eta_s^2(\mu_{ij}^2+\sigma_{ij}^2)$.

\textbf{Step 3: When skipping helps.}
Skipping entry $(i,j)$ is beneficial when $\Delta R_{ij} \leq 0$, i.e.,
$2|\mu_{ij}||e_{ij}| \leq \eta_s(\mu_{ij}^2+\sigma_{ij}^2)$.
This is the state-dependent condition stated in the proposition.
Summing over all skipped entries and $T/B$ periods gives the bound.
\end{proof}

\begin{remark}[When selective consolidation can fail]
\label{rem:selective-fail}
If the threshold $\tau_s$ is set too aggressively (too high), signal
coordinates may be misclassified as transient, causing the learner to
miss genuine subspace drift. The synthetic experiment E4 showed this
failure mode with the previous covariance-threshold criterion, which
was too sensitive to the noise distribution.
The gradient-sign criterion is more robust because it directly measures
the consistency of the optimization direction.
\end{remark}
